% Bab I  Pendahuluan
\chapter{PENDAHULUAN}

\section{Latar Belakang}
Latar belakang berisi justifikasi terhadap penelitian yang dilakukan: permasalahan yang dihadapi secara umum, kesenjangan antara kondisi yang diharapkan dan kondisi nyata, serta alasan penelitian ini penting dan layak dilakukan. Dukung uraian dengan data dan pendapat peneliti lain yang relevan, dengan sitasi gaya APA, misalnya Andrews (2007) menyatakan bahwa ... atau ... (Capozzi et al., 2013).

\textit{Contoh kalimat bersitasi:} Penjadwalan paket pada jaringan seluler berperan penting dalam membagi sumber daya radio secara adil di antara pengguna (Capozzi et al., 2013). Andrews (2007) menyatakan bahwa teori penjadwalan pada jaringan data nirkabel terus berkembang seiring bertambahnya kebutuhan layanan. Kebutuhan serupa juga muncul pada jaringan listrik cerdas yang menuntut penjadwalan waktu nyata (Xu \& Fischione, 2012). Data statistik nasional dipublikasikan secara berkala oleh lembaga resmi (Badan Pusat Statistik, 1997).

\section{Rumusan Masalah}
Rumusan masalah disusun berdasarkan latar belakang, berbentuk pertanyaan penelitian atau pernyataan yang memerlukan jawaban melalui penelitian ilmiah:
\begin{enumerate}
\item Bagaimana ... ?
\item Seberapa besar ... ?
\end{enumerate}

\section{Tujuan Penelitian}
Tujuan penelitian dinyatakan dalam kalimat singkat, jelas, dan spesifik, serta konsisten dengan rumusan masalah:
\begin{enumerate}
\item Merancang ...
\item Mengukur ...
\end{enumerate}

\section{Manfaat Penelitian}
Uraikan manfaat yang diperoleh apabila tujuan tercapai, difokuskan pada pengembangan keilmuan, informasi bagi pembuat kebijakan, dan umpan balik bagi pengguna.

\section{Batasan Masalah}
Tuliskan batasan agar penelitian terfokus, misalnya ruang lingkup objek, perangkat, atau data yang digunakan.

\textit{Contoh kalimat bersitasi:} Perangkat keras yang digunakan mengikuti spesifikasi dari pabrikan (Espressif Systems, 2023).

\section{Sistematika Penulisan}
Skripsi ini terdiri atas lima bab. Bab I berisi pendahuluan. Bab II berisi tinjauan pustaka. Bab III berisi metode penelitian. Bab IV berisi hasil penelitian dan pembahasan. Bab V berisi kesimpulan dan saran.

